% New analytic audit, 21 October 2026. Insert after the existing Appendices A--C.
% Uses the article's macros \X and \nrm. No new theorem counter is introduced.
\section{The full dual residual of the neutral reference}
\label{app:reflected-residual}

This calculation justifies the exterior residual formula used in the
fixed critical-state certificate. In centered coordinates
$(\xi,y)=(x_1-6,x_2)$, set $a=L/2=6$ and
\[
 \varphi(\xi,y)=\arg(\xi+a+\mathrm i y)
                  -\arg(\xi-a+\mathrm i y),\qquad L=12.
\]
Use the real branch tending to zero away from the finite cut
$[-a,a]\times\{0\}$. Its gradient is square summable; compact cutoffs
converge in the energy norm since $\varphi=O(L/r)$ and
$D\varphi=O(L/r^2)$ outside that cut. Thus $\varphi\in\X$.
Each unoriented bond is counted once, with
$D_eu=u(e^+)-u(e^-)$, and $D^*$ denotes the vertex-distribution adjoint.
Write $f_\theta=D^*\sin(D\theta)$.

Let $S>a+2$ be an integer and define the symmetric vertex square
\[
 B_S=\{v\in(\mathbb Z+1/2)^2:\nrm{v}_\infty\le S-3/2\},
 \qquad b=2S-2,\qquad \chi=\frac{S}{S-a-2}.
\]
Here the square is in the centered coordinates. Suppose $h$ is compact
and every bond on which $D_eh\ne0$ has \emph{both} endpoints in $B_S$.
Set $f_B=\mathbf1_{B_S}f_{\varphi+h}$, interpreted as a load on the
full lattice. Then $f_B$ has algebraically zero mean. Indeed,
$j_e=\sin D_e(\varphi+h)-\sin D_e\varphi$ is a compact flux whose
nonzero bonds and both endpoints lie in $B_S$, so
$\sum_{B_S}D^*j=0$ by endpoint cancellation. The pure-reference phase is
complex-conjugated by $(\xi,y)\mapsto(\xi,-y)$; its vertex force is
therefore odd under this reflection. Pairing the vertices of $B_S$
gives $\sum_{B_S}f_\varphi=0$. No oddness assumption on $h$ is needed.
An interval sum containing zero would not replace this argument.

Let $L_N$ and $D_N$ count only bonds interior to $B_S$. For any finite
proposal $u$, the Neumann comparison and
$\lambda_1(L_N)=4\sin^2(\pi/(2b))\ge4/b^2$ give
\begin{equation}
 \nrm{f_B}_{\X^*}
 \le \nrm{D_Nu}_2+\frac b2\nrm{f_B-L_Nu}_2.
 \label{eq:finite-residual-energy}
\end{equation}
To see this, solve $L_Nu_*=f_B$ with the same mean as $u$, and use
$\nrm{f_B}_{\X^*}\le\nrm{D_Nu_*}_2$ and
$\nrm{D_N(u_*-u)}_2\le\lambda_1(L_N)^{-1/2}
\nrm{f_B-L_Nu}_2$. No normalization of the mean of $u$ is required.

The complementary load obeys the explicit bound
\begin{equation}
 \nrm{f_{\varphi+h}-f_B}_{\X^*}
 \le \frac{64\sqrt3}{5}\frac{L\chi^5}{(S-1)^{5/2}}
      +\frac{32\sqrt3}{9}\frac{L^3\chi^7}{(S-1)^{9/2}}.
 \label{eq:reflected-residual-tail}
\end{equation}
Here is a derivation including the cancellation needed for the second
term. In a differentiation region at distance at least $d$ from the cut,
subtracting the complex derivatives of the two argument fields gives
\[
 |\partial_i^r\varphi|\le r!L/d^{r+1}\quad(r\ge1),
 \qquad \Delta\varphi=0.
\]
For $p_i=\varphi(v+e_i)-\varphi(v)$ and
$m_i=\varphi(v)-\varphi(v-e_i)$, the central Taylor remainder of order
four and harmonicity yield $|L_\infty\varphi(v)|\le4L/d^5$.
Also $|p_i|,|m_i|\le L/d^2$ and $|p_i-m_i|\le2L/d^3$.
Since $g(t)=\sin t-t$ has $|g'(t)|\le t^2/2$,
\[
 |g(p_i)-g(m_i)|\le L^3/d^7,\qquad
 |f_\varphi(v)|\le4L/d^5+2L^3/d^7.
\]
The two opposite bonds are combined before taking absolute values.

Exterior shells have $\nrm{v}_\infty=k-1/2$, $k\ge S$, and at most
$8k$ vertices. The unit segments used above lie at distance at least
$k-a-3/2\ge k-a-2$ from the cut, while
$(k-a-2)^{-r}\le\chi^r k^{-r}$. Pair each exterior vertex with $y>0$
with its reflected vertex $Rv$. The load $\delta_v-\delta_{Rv}$ has
$\X^*$ norm at most $\sqrt{2|y|}\le\sqrt{3k}$, by Cauchy--Schwarz
along the vertical lattice path. This path may enter $B_S$; only its
endpoints specify the exterior load. Since $h$ changes no exterior
vertex force, the triangle inequality gives
\begin{align*}
 \nrm{f_{\varphi+h}-f_B}_{\X^*}
 &\le32\sqrt3 L\chi^5\sum_{k\ge S}k^{-7/2}
       +16\sqrt3 L^3\chi^7\sum_{k\ge S}k^{-11/2}.
\end{align*}
Integrating these positive decreasing terms from $S-1$ proves
\eqref{eq:reflected-residual-tail}. The convergent majorant justifies
the load series in $\X^*$; equality with the exterior vertex force is
first checked on compact tests and then extended by density.

In the checker the correction array has even side $n$, its boundary
is zero, the phase box has side $n+4$, and the force box has side
$b=n+2$. Its halo includes every changed flux bond and both endpoints,
and all four neighbours of every force-box vertex are present.
Thus $S=b/2+1$ and $\chi=S/(S-8)$ agree exactly with the implemented
formula. Adding \eqref{eq:finite-residual-energy} and
\eqref{eq:reflected-residual-tail} bounds the full dual residual.
These estimates concern the fixed neutral reference and its supported
corrections; they do not by themselves establish a barrier or a global
separator.
